% Generated draft: requires mathematical and source review.
\documentclass[11pt]{article}
\usepackage[margin=25mm]{geometry}
\usepackage{fontspec}
\setmainfont{Latin Modern Roman}
\usepackage{amsmath,amssymb,mathtools,unicode-math}
\setmathfont{Latin Modern Math}
\usepackage{enumitem,needspace,xurl,hyperref}
\setlength{\parindent}{0pt}
\setlength{\parskip}{5pt}
\setlength{\emergencystretch}{3em}
\newcommand{\N}{\mathbb{N}}
\newcommand{\Z}{\mathbb{Z}}
\newcommand{\Q}{\mathbb{Q}}
\newcommand{\R}{\mathbb{R}}
\newcommand{\C}{\mathbb{C}}
\newcommand{\F}{\mathbb{F}}
\newcommand{\A}{\mathbb{A}}
\newcommand{\PP}{\mathbb{P}}
\newcommand{\E}{\mathbb{E}}
\newcommand{\eps}{\varepsilon}
\newcommand{\dd}{\,\mathrm{d}}
\begin{document}
\section*{UnsolvedMath ID: 30005800}
\textbf{Source problem code:} catalogue entry
\subsection{Definitions and mathematical statement}
For a space \(X\) and a coefficient field \(K\), write \(\widetilde\beta_i(X;K)=\dim_K\widetilde H_i(X;K)\) for its reduced singular Betti number. The supplied record does not specify the coefficient field or homology theory; this notation explains a customary interpretation without resolving that omission.

The supplied conjecture fixes constants \(c,t\) and considers a family \(\mathcal F\) of subsets of \(\R^d\) satisfying
\[
\sum_{i=1}^{\lfloor d/2\rfloor}\widetilde\beta_i\!\left(\bigcap_{S\in\mathcal G}S;K\right)\leq c\,|\mathcal G|^t
\qquad\text{for every finite }\mathcal G\subseteq\mathcal F.
\]
Here \(d\) is the ambient Euclidean dimension. The permitted ranges of \(c,t\), conventions for \(\mathcal G=\varnothing\), and any regularity assumptions on the sets are not supplied. In particular, the displayed sum starts at \(i=1\); it does not include \(\widetilde\beta_0\). When \(\lfloor d/2\rfloor=0\), the sum is empty and equals zero.

For integers \(p\geq q\geq1\), a finite family \(\mathcal H\) has the \((p,q)\)-property if every \(p\)-element subfamily \(\mathcal A\subseteq\mathcal H\) contains a \(q\)-element subfamily \(\mathcal B\subseteq\mathcal A\) with
\[
\bigcap_{S\in\mathcal B}S\neq\varnothing.
\]
Its piercing number is
\[
\tau(\mathcal H)=\min\left\{|P|:P\subseteq\R^d\text{ is finite and }P\cap S\neq\varnothing\text{ for every }S\in\mathcal H\right\},
\]
with value \(+\infty\) if no such finite set exists.

\textbf{Requested conjectural conclusion.} Such families satisfy a \((p,q)\)-theorem for every pair of integers \(p\geq q\geq d+1\). In the customary uniform finite-family interpretation, this asks for a finite bound \(N=N(d,c,t,p,q)\), independent of \(\mathcal F\) and \(\mathcal H\), such that every finite \(\mathcal H\subseteq\mathcal F\) with the \((p,q)\)-property satisfies \(\tau(\mathcal H)\leq N\). The homology coefficients must also be fixed when interpreting this assertion.

\textbf{Unresolved qualifications.} The supplied record does not explicitly define the intended \((p,q)\)-theorem, its uniform dependence, or whether it concerns finite subfamilies or a single transversal for an infinite family. It also does not exclude empty members or specify how families with fewer than \(p\) members are treated; for such families the displayed \((p,q)\)-property is vacuous. These conventions require clarification. The absence of a degree-zero homology bound, especially in dimensions where the sum is empty, requires checking against the intended source hypotheses. No connectedness assumption, degree-zero term, or other missing hypothesis has been inserted.
\subsection{Short English statement}
The supplied conjecture asks whether polynomial bounds on the indicated reduced Betti numbers of all finite intersections imply a uniform bound on piercing numbers whenever the \((p,q)\)-property holds, for every \(p\geq q\geq d+1\). The record labels the problem open. Homology conventions, possible missing hypotheses, and the precise finite or infinite piercing formulation remain unresolved; this is a statement-only draft requiring mathematical review.
\subsection{Sources}
\begin{itemize}
\item Local UnsolvedMath registry: supplied source record, permanent ID \(30005800\), title \emph{Polynomial-Homology Families and the p-q Theorem}, exact source problem code \texttt{catalogue entry}; collection: \emph{Oberwolfach Reports Open Problems}. The local registry files were not inspected for this draft.
\item Supplied report citation and literature-source link: \emph{Discrete Geometry} (2024), Oberwolfach Reports, DOI \texttt{10.4171/owr/2024/3}, pp.\ 137--202, \url{https://doi.org/10.4171/owr/2024/3}. No conjecture-specific page, contribution title, or proposer is supplied.
\item Supplied background attributes Theorem 7 to Goaoc--Holmsen--Pat\'akov\'a (2021): if \(\mathrm{HC}_{\lfloor d/2\rfloor}(\mathcal F)\leq C\), then the fractional Helly number is at most \(d+1\), with the stated consequence of a \((p,q)\)-theorem for \(p\geq q\geq d+1\). The record does not define \(\mathrm{HC}\) or provide a separate bibliographic reference. These are reported background claims, not independently verified assertions in this draft.
\item The supplied background cites earlier conjectures of Kalai and Meshulam as motivation and mentions examples from geometric transversal theory suggesting that constant homological complexity need not be necessary. No precise references or examples are supplied.
\item Literature assessment dated exactly \texttt{2026-08-22}: supplied dated curation labels the problem open and reports that the 2024 report presents the assertion as a conjecture, while the cited fractional-Helly result covers a special regime rather than the full assertion. This assessment is not independently verified. The supplied record also labels the statement \texttt{corrected\_verified} and reports reconstruction and checking against the official report; that verification is source metadata, not verification performed for this draft.
\end{itemize}
\end{document}
